\honest{Thou Art Godshatter}{Eliezer Yudkowsky}{2007}

I open with a puzzle. Before the twentieth century nobody had the concept of ``inclusive genetic fitness,'' the one thing natural selection favours, and our brains grant sexual pleasure without checking whether sex can produce children. Why did the ``Evolution-of-Humans Fairy'' build brains that would invent condoms?

A human designer might think it easy to make agents that want nothing but fitness. They would have sex only to reproduce, eat only because food serves reproduction, and, as grandmothers, kill themselves once they cost more than they help. I answer two objections to show that such an agent needs no separate drives. If it knows that it must eat to reproduce, it will eat. If it knows that a kind of curiosity paid off for its ancestors, it will be curious. So the pure fitness maximizer is possible ``in principle.''

Evolution did not build one because ``the blind idiot god isn't that smart'': it cannot rewrite a whole design at once, as a programmer can. For millions of years brains learned from rewards, such as the taste of fat and the pleasure of sex, that went with reproduction only at long range. When reasoning appeared, it was put to work getting those rewards, ``a relatively simple hack.'' \nb{This is an inference from what selection could easily find. The post states it as history and gives no evidence about hominid evolution.} A human engineer would have made fitness explicit, I add, ``But then a human engineer wouldn't have built the retina backward, either.'' \nb{The retina example is set out in ``An Alien God,'' earlier in the book. Some read the retina's layout as a functional compromise, not only a relic of history.}

So evolution's single goal did not pass into us. Its ``optimization criterion did not successfully quine,'' and it splintered ``into a thousand shards of desire.'' Maybe, given a billion more years, selection would produce true fitness maximizers, as in Niven and Pournelle's \textsc{The Mote in God's Eye}. \nb{In Wikipedia's summary of the novel, its aliens must become pregnant periodically or die. That is a compulsion to breed; whether the novel presents them as fitness maximizers could not be checked.}

We love sugar, our children, status, sex, play and learning, whether or not they help us reproduce. ``And this is well, I think, though I'm a human who says so.'' The alternative would be descendants who fill the galaxies with replicators. Our tastes for novelty and challenge find evolution's single aim unsatisfying. We got those tastes from evolution too, from ``the blind idiot's godshatter.'' \nb{The word comes from Vernor Vinge's novel \textsc{A Fire Upon the Deep} (1992). The post does not say so.} ``So what?'' \nb{That states a position, that where a value came from does not settle whether to keep it. The post gives no argument for it.}
